% GENERATED FILE. Edit canonical lessons, then run npm run generate:books.
% canonical-source-hash: fnv1a64:9436ef40065534ad
% canonical-lessons: TE-C177-okavela, TE-C177-badulu, TE-C177-prakaram, TE-C177-asalu, TE-C177-edemaina

\chapter{If by chance, Instead of, According to, Actually, Anyway}
\label{ch:te-if-by-chance-instead-of-according-to-actually-anyway}
\hlchaptermodality{177}

\begin{chapteropening}
\textbf{By the end of this chapter:} \emph{I can say if by chance, instead of, according to, actually and anyway.}

Complete the last lesson of chapter 177: I can say if by chance, instead of, according to, actually and anyway.
\end{chapteropening}

\section[\texorpdfstring{\te{ఒకవేళ} (okavēḷa)}{ఒకవేళ (okavēḷa)}]{\te{ఒకవేళ} (okavēḷa) — if by chance, in case}

\label{lesson:TE-C177-okavela}



\begin{quote}
Before the new one: say the Telugu for about, then the Telugu for without.
\end{quote}

\subsection*{You'll want to know: \te{ఒకవేళ}}
\textbf{\te{ఒకవేళ}} — \emph{okavēḷa} — \textquotedblleft{}if by chance, in case\textquotedblright{}.

\begin{grammarlens}[title={What you've built}]
One more word for this chapter.
\end{grammarlens}

\subsection*{Your turn}
\begin{itemize}
\raggedright
  \item \emph{Say it:} \emph{okavēḷa}
  \item \emph{Say it:} \emph{okavēḷa}, once more
  \item \emph{Say it:} say \emph{okavēḷa}
  \item \emph{Recall:} say the Telugu for about, then the Telugu for without, then say \emph{okavēḷa} again
\end{itemize}

\begin{tcolorbox}[breakable,colback=teal!4,colframe=teal!35!black,title={Before you move on}]
What does \te{ఒకవేళ} mean? (\textquotedblleft{}If by chance, in case\textquotedblright{}.) Say it once more.
\end{tcolorbox}
\section[\texorpdfstring{\te{బదులు} (badulu)}{బదులు (badulu)}]{\te{బదులు} (badulu) — instead of}

\label{lesson:TE-C177-badulu}



\begin{quote}
Before the new one: say the Telugu for without, then the Telugu for if by chance.
\end{quote}

\subsection*{You'll want to know: \te{బదులు}}
\textbf{\te{బదులు}} — \emph{badulu} — \textquotedblleft{}instead of\textquotedblright{}.

\begin{grammarlens}[title={What you've built}]
One more word for this chapter.
\end{grammarlens}

\subsection*{Your turn}
\begin{itemize}
\raggedright
  \item \emph{Say it:} \emph{badulu}
  \item \emph{Say it:} \emph{badulu}, once more
  \item \emph{Say it:} say \emph{badulu}
  \item \emph{Recall:} say the Telugu for without, then the Telugu for if by chance, then say \emph{badulu} again
\end{itemize}

\begin{tcolorbox}[breakable,colback=teal!4,colframe=teal!35!black,title={Before you move on}]
What does \te{బదులు} mean? (\textquotedblleft{}Instead of\textquotedblright{}.) Say it once more.
\end{tcolorbox}
\section[\texorpdfstring{\te{ప్రకారం} (prakāraṁ)}{ప్రకారం (prakāraṁ)}]{\te{ప్రకారం} (prakāraṁ) — according to}

\label{lesson:TE-C177-prakaram}



\begin{quote}
Before the new one: say the Telugu for if by chance, then the Telugu for instead of.
\end{quote}

\subsection*{You'll want to know: \te{ప్రకారం}}
\textbf{\te{ప్రకారం}} — \emph{prakāraṁ} — \textquotedblleft{}according to\textquotedblright{}.

\begin{grammarlens}[title={What you've built}]
One more word for this chapter.
\end{grammarlens}

\subsection*{Your turn}
\begin{itemize}
\raggedright
  \item \emph{Say it:} \emph{prakāraṁ}
  \item \emph{Say it:} \emph{prakāraṁ}, once more
  \item \emph{Say it:} say \emph{prakāraṁ}
  \item \emph{Recall:} say the Telugu for if by chance, then the Telugu for instead of, then say \emph{prakāraṁ} again
\end{itemize}

\begin{tcolorbox}[breakable,colback=teal!4,colframe=teal!35!black,title={Before you move on}]
What does \te{ప్రకారం} mean? (\textquotedblleft{}According to\textquotedblright{}.) Say it once more.
\end{tcolorbox}
\section[\texorpdfstring{\te{అసలు} (asalu)}{అసలు (asalu)}]{\te{అసలు} (asalu) — actually; (not) at all}

\label{lesson:TE-C177-asalu}



\begin{quote}
Before the new one: say the Telugu for instead of, then the Telugu for according to.
\end{quote}

\subsection*{You'll want to know: \te{అసలు}}
\textbf{\te{అసలు}} — \emph{asalu} — \textquotedblleft{}actually; (not) at all\textquotedblright{}.

\begin{grammarlens}[title={What you've built}]
One more word for this chapter.
\end{grammarlens}

\subsection*{Your turn}
\begin{itemize}
\raggedright
  \item \emph{Say it:} \emph{asalu}
  \item \emph{Say it:} \emph{asalu}, once more
  \item \emph{Say it:} say \emph{asalu}
  \item \emph{Recall:} say the Telugu for instead of, then the Telugu for according to, then say \emph{asalu} again
\end{itemize}

\begin{tcolorbox}[breakable,colback=teal!4,colframe=teal!35!black,title={Before you move on}]
What does \te{అసలు} mean? (\textquotedblleft{}Actually; (not) at all\textquotedblright{}.) Say it once more.
\end{tcolorbox}
\section[\texorpdfstring{\te{ఏదేమైనా} (ēdēmainā)}{ఏదేమైనా (ēdēmainā)}]{\te{ఏదేమైనా} (ēdēmainā) — anyway, in any case}

\label{lesson:TE-C177-edemaina}



\begin{quote}
Before the new one: say the Telugu for according to, then the Telugu for actually.
\end{quote}

\subsection*{You'll want to know: \te{ఏదేమైనా}}
\textbf{\te{ఏదేమైనా}} — \emph{ēdēmainā} — \textquotedblleft{}anyway, in any case\textquotedblright{}.

\begin{grammarlens}[title={What you've built}]
One more word for this chapter.
\end{grammarlens}

\subsection*{Your turn}
\begin{itemize}
\raggedright
  \item \emph{Say it:} \emph{ēdēmainā}
  \item \emph{Say it:} \emph{ēdēmainā}, once more
  \item \emph{Say it:} say \emph{ēdēmainā}
  \item \emph{Recall:} say the Telugu for according to, then the Telugu for actually, then say \emph{ēdēmainā} again
\end{itemize}

\begin{tcolorbox}[breakable,colback=teal!4,colframe=teal!35!black,title={Before you move on}]
What does \te{ఏదేమైనా} mean? (\textquotedblleft{}Anyway, in any case\textquotedblright{}.) Say it once more.
\end{tcolorbox}
